\section{O \emph{benchmark} de Hysing: bolha ascendente}
\label{sec:hysing}

\subsection{Por que este teste, e por que ele \'e o \'ultimo}

Todos os or\'aculos anteriores foram \emph{internos} --- anal\'iticos (Laplace,
Lamb, Rider-Kothe) ou m\'utuos (um m\'etodo contra o outro). S\~ao fortes, e
foram o que localizou cada defeito deste relat\'orio, mas t\^em um limite
comum: nenhum deles diz se o \emph{par} de implementa\c{c}\~oes concorda com o
que outros grupos obt\^em no mesmo problema. Duas implementa\c{c}\~oes que
compartilham um erro de formula\c{c}\~ao concordariam entre si perfeitamente.

Hysing et al.~\cite{hysing2009} fecham essa lacuna. O artigo publica, para um
problema de bolha ascendente completamente especificado, os valores obtidos por
\textbf{tr\^es c\'odigos independentes} (TP2D, FreeLIFE e MooNMD), com
resolu\c{c}\~ao refinada at\'e a converg\^encia e com as discrep\^ancias entre
eles declaradas. \'E um or\'aculo \emph{externo}, e com barra de erro pr\'opria.

\paragraph{O caso 1, exatamente como especificado.} Dom\'inio $[0,1]\times[0,2]$,
bolha circular de raio $R=0{,}25$ centrada em $(0{,}5;\,0{,}5)$, sem escorregamento
em cima e embaixo e paredes laterais, at\'e $t=3$:

\begin{center}
\begin{tabular}{lrr}
\toprule
 & fase externa (1) & bolha (2) \\
\midrule
densidade $\rho$    & 1000 & 100 \\
viscosidade $\mu$   & 10   & 1   \\
\bottomrule
\end{tabular}
\qquad
\begin{tabular}{lr}
\toprule
$\sigma$ & 24{,}5 \\
$g$      & 0{,}98 \\
$Re$     & 35 \\
$E\ddot{o}$ & 10 \\
\bottomrule
\end{tabular}
\end{center}

Nesse regime a bolha se deforma mas \textbf{n\~ao se parte} --- condi\c{c}\~ao
para que o front-tracking possa concorrer, j\'a que mudan\c{c}a de topologia est\'a
fora do escopo deste trabalho (Se\c{c}\~ao~\ref{sec:discussao}). Foi por isso
que se escolheu o caso 1 e n\~ao o caso 2, em que a bolha desenvolve filamentos
e se rompe.

\subsection{Os valores publicados}

\pgfplotstabletypeset[
  string type,
  columns/grandeza/.style={column name={grandeza}},
  columns/TP2D/.style={column name={TP2D}},
  columns/FreeLIFE/.style={column name={FreeLIFE}},
  columns/MooNMD/.style={column name={MooNMD}},
  columns/consenso/.style={column name={consenso}},
  every head row/.style={before row=\toprule, after row=\midrule},
  every last row/.style={after row=\bottomrule},
]{\dat{hysing-referencia.dat}}

A coluna \emph{consenso} \'e o que serve de alvo: onde os tr\^es c\'odigos
concordam, a faixa \'e estreita (a circularidade m\'inima em
$0{,}9012\pm0{,}0001$) e uma discrep\^ancia nossa acima dela \'e defeito nosso;
onde eles divergem (o centro de massa em $t=3$ varia de $1{,}0799$ a $1{,}0817$
entre os tr\^es), a faixa \'e o pr\'oprio limite do or\'aculo, e cair dentro
dela \'e tudo o que se pode exigir.

\subsection{Uma defini\c{c}\~ao errada, e o erro que ela inventou}

Antes dos resultados, um erro meu que vale registrar porque n\~ao se manifestou
como falha, e sim como um n\'umero ruim plaus\'ivel.

A circularidade de Hysing \'e a raz\~ao entre o per\'imetro do c\'irculo de
\'area equivalente e o per\'imetro da bolha:
\begin{equation}
  c = \frac{P_a}{P_b} = \frac{2\sqrt{\pi A}}{P} .
  \label{eq:circ-hysing}
\end{equation}
Eu vinha medindo $4\pi A/P^{2}$, que \'e o \textbf{quadrado} de
\eqref{eq:circ-hysing}. As duas s\~ao 1 no c\'irculo, as duas decrescem com a
deforma\c{c}\~ao, e as duas ficam em $[0,1]$ --- nada no comportamento denuncia a
troca. Mas elevar ao quadrado um n\'umero perto de $0{,}9$ afasta-o do alvo:
medindo pela minha, reportei \SI{11}{\percent} de erro onde o valor compar\'avel
ao do artigo era \SI{0.83}{\percent}. O erro estava na defini\c{c}\~ao, e n\~ao
na simula\c{c}\~ao.

A li\c{c}\~ao \'e espec\'ifica e \'util: comparar com valor publicado exige ler a
\emph{defini\c{c}\~ao} publicada, e n\~ao reconhecer a grandeza pelo nome. Duas
f\'ormulas que se comportam igual em todos os casos-limite podem diferir por
composi\c{c}\~ao com uma fun\c{c}\~ao mon\'otona, e a compara\c{c}\~ao com um
n\'umero de fora \'e justamente onde isso importa.

\paragraph{A circularidade s\'o \'e reportada do lado front-tracking.} Ela pede o
per\'imetro da interface. No front-tracking ele sai em forma fechada do
pol\'igono; no VOF exigiria a reconstru\c{c}\~ao PLIC, que o solver n\~ao
exp\~oe. Preferiu-se \textbf{faltar} a aproximar por outro caminho: com duas
aproxima\c{c}\~oes diferentes, a compara\c{c}\~ao entre os m\'etodos viraria
compara\c{c}\~ao entre defini\c{c}\~oes --- exatamente o erro do par\'agrafo
anterior, agora entre as duas colunas em vez de contra o artigo. O centro de
massa $y_c$ e a velocidade de ascens\~ao $V_c$ saem da integral do campo de
fra\c{c}\~ao, que os dois m\'etodos t\^em, pela \emph{mesma} rotina
(\texttt{examples-common/hysing-metricas.c}).

\subsection{O que este caso exigiu de maquinaria}

Vale listar, porque a lista \'e o conte\'udo real das se\c{c}\~oes anteriores e
porque cada item foi um defeito encontrado:

\begin{itemize}
  \item \textbf{Malha adaptativa nos dois m\'etodos}, pelo mesmo crit\'erio e
    pela m\'aquina verificada de remalhamento, com a regra das cinco c\'elulas
    finas por lado cumprida e medida (Se\c{c}\~ao~\ref{sec:amr}).
  \item \textbf{Espa\c{c}amento de marcadores derivado da malha}, e n\~ao fixado
    \`a m\~ao: com 64 marcadores em $h=1/160$ a frente perde \num{2.1e-2} de
    \'area; derivando $\Delta s \approx h$, \num{1.95e-5}. Aqui s\~ao 252
    marcadores, de $\Delta s/h = 1{,}00$.
  \item \textbf{Restri\c{c}\~ao capilar de passo.} $\Delta t < \sqrt{\rho
    h^{3}/2\pi\sigma}$ escala como $h^{3/2}$, e com dois n\'iveis de refino
    aperta 8 vezes; daqui vem $\Delta t = \num{2.5e-4}$, e as 12\,000 itera\c{c}\~oes
    que $t=3$ custa.
  \item \textbf{Persist\^encia da frente.} Uma corrida de 12\,000 passos leva
    cerca de 6~h nesta m\'aquina, e a primeira tentativa morreu no meio. O
    \texttt{h.save} do solver guarda \emph{campos}; a frente lagrangeana n\~ao \'e
    campo, e nenhum campo a determina. Retomar sem grav\'a-la poria os campos no
    instante salvo e os marcadores em $t=0$ --- e rodaria, em sil\^encio. A
    grava\c{c}\~ao foi implementada e verificada: advectar 100 passos direto e
    advectar 50, gravar, ler e advectar os outros 50 d\~ao o mesmo resultado
    \textbf{bit a bit}. Do lado VOF, retomando do passo 1401 da corrida
    interrompida, 19 dos 23 passos em comum saem id\^enticos em todos os d\'igitos
    impressos e 4 diferem no \'ultimo --- \num{1e-10}, a toler\^ancia do solver
    linear.
\end{itemize}

\subsection{Resultado parcial: os dois m\'etodos concordam entre si}

\paragraph{Esta subse\c{c}\~ao est\'a incompleta, e de prop\'osito.} As corridas
at\'e $t=3$ estavam em andamento quando este texto foi escrito. O fechamento
contra os valores publicados --- circularidade m\'inima, velocidade m\'axima e
centro de massa em $t=3$ --- \textbf{n\~ao est\'a medido aqui}, e a
Se\c{c}\~ao~\ref{sec:amr} e a lista de pr\'oximos passos continuam valendo quando
dizem que o Hysing segue sem fechar. O que est\'a abaixo \'e o trecho
$t \in [0;\,0{,}35]$, medido nos dois m\'etodos antes da interrup\c{c}\~ao.

\pgfplotstabletypeset[
  string type,
  columns/t/.style={column name={$t$}},
  columns/ycft/.style={column name={$y_c$ FT}},
  columns/ycvof/.style={column name={$y_c$ VOF}},
  columns/difyc/.style={column name={dif.}},
  columns/vcft/.style={column name={$V_c$ FT}},
  columns/vcvof/.style={column name={$V_c$ VOF}},
  columns/difvc/.style={column name={dif.}},
  columns/circft/.style={column name={$c$ FT}},
  every head row/.style={before row=\toprule, after row=\midrule},
  every last row/.style={after row=\bottomrule},
]{\dat{hysing-parcial.dat}}

\begin{figure}[htbp]
\centering
\begin{tikzpicture}
\begin{axis}[
  name=cheia, width=6.4cm, height=6.4cm, axis equal image,
  xlabel={$x$}, ylabel={$y$}, grid=major,
  xmin=0.23, xmax=0.77, ymin=0.28, ymax=0.82,
  title={$t=0{,}45$, a bolha inteira}, title style={font=\small},
  legend style={font=\scriptsize, at={(0.5,-0.22)}, anchor=north},
  legend columns=2, legend cell align=left,
]
\addplot[thick]         table[x=x,y=y] {\dat{formas/hy_ft_3.dat}};
\addlegendentry{front-tracking}
\addplot[densely dashed] table[x=x,y=y] {\dat{formas/hy_vof_3.dat}};
\addlegendentry{VOF}
\end{axis}
\begin{axis}[
  at={(cheia.right of south east)}, anchor=left of south west,
  width=6.4cm, height=6.4cm, axis equal image,
  xlabel={$x$}, grid=major,
  xmin=0.676, xmax=0.692, ymin=0.365, ymax=0.381,
  title={ampliado onde a diferen\c{c}a \'e maior}, title style={font=\small},
  scaled ticks=false, tick label style={font=\tiny, /pgf/number format/fixed},
]
\addplot[thick]         table[x=x,y=y] {\dat{formas/hy_ft_3.dat}};
\addplot[densely dashed] table[x=x,y=y] {\dat{formas/hy_vof_3.dat}};
\end{axis}
\end{tikzpicture}
\caption{A interface nos dois m\'etodos no mesmo instante. \`A esquerda, na
escala da bolha, as duas curvas s\~ao indistingu\'iveis; \`a direita, ampliado
25 vezes em torno do ponto de maior discrep\^ancia, aparece a separa\c{c}\~ao de
\num{1.7e-4} --- cerca de \SI{2.7}{\percent} de uma c\'elula fina. \textbf{Em
$t=0{,}45$ a bolha ainda \'e praticamente circular}, ent\~ao esta figura mostra
a \emph{concord\^ancia} entre as duas representa\c{c}\~oes, e n\~ao
deforma\c{c}\~ao: a forma deformada, que \'e o que este \emph{benchmark} testa,
s\'o aparece a partir de $t\approx1$ e est\'a fora do alcance dos dados
dispon\'iveis para este texto. A curva do VOF \'e o contorno
$\mathrm{FracVol}=0{,}5$ obtido por \emph{marching squares} com os valores de
c\'elula medianizados nos v\'ertices --- confer\~ido contra o c\'irculo
anal\'itico em $t=0$ (desvio radial m\'aximo \num{1.34e-4}), e adequado para
desenhar a forma, n\~ao para medir per\'imetro.}
\label{fig:hysing-forma}
\end{figure}

A concord\^ancia de forma, medida:

\pgfplotstabletypeset[
  string type,
  columns/t/.style={column name={$t$}},
  columns/areaft/.style={column name={\'area FT}},
  columns/areavof/.style={column name={\'area VOF}},
  columns/difarea/.style={column name={dif.}},
  columns/desloccentro/.style={column name={$|\Delta y_c|$}},
  columns/difradial/.style={column name={dif. radial m\'ax.}},
  columns/emdoraio/.style={column name={\% do raio}},
  every head row/.style={before row=\toprule, after row=\midrule},
  every last row/.style={after row=\bottomrule},
]{\dat{hysing-forma.dat}}

A leitura desta tabela tem uma sutileza que convém n\~ao perder: das
\num{1.65e-4} de discrep\^ancia radial em $t=0{,}45$, \num{1.06e-4} j\'a existem
em $t=0$ --- s\~ao a diferen\c{c}a de \emph{inicializa\c{c}\~ao} entre um
pol\'igono de 252 lados inscrito no c\'irculo e o contorno extra\'ido do campo de
fra\c{c}\~ao, e n\~ao diverg\^encia din\^amica. O que as duas
discretiza\c{c}\~oes acumulam de fato em 1800 passos \'e a diferen\c{c}a entre
esses dois n\'umeros, cerca de \num{6e-5}. \'E o mesmo efeito que a
Se\c{c}\~ao~\ref{sec:comparacao} registrou no d\'eficit de \'area do pol\'igono:
uma parte do que parece erro de m\'etodo \'e escolha de inicializa\c{c}\~ao.

Duas leituras, e nenhuma delas \'e o fechamento.

\paragraph{1. A concord\^ancia m\'utua \'e forte.} O centro de massa difere em
no m\'aximo \num{1e-5} e a velocidade de ascens\~ao em \num{7.4e-5} ao longo de
1400 passos --- quatro a cinco ordens de grandeza abaixo das pr\'oprias
grandezas. S\~ao duas representa\c{c}\~oes de interface completamente diferentes,
uma lagrangeana e uma euleriana, com c\'alculos de curvatura e de for\c{c}a sem
nada em comum, produzindo a mesma trajet\'oria. Isso \emph{n\~ao} prova que a
trajet\'oria est\'a certa --- \'e a limita\c{c}\~ao do or\'aculo m\'utuo, e \'e
exatamente por isso que o or\'aculo externo importa --- mas torna improv\'avel
que o erro esteja em algo que os dois m\'etodos n\~ao compartilham, e concentra
a suspeita no que \'e comum: malha, passo, condi\c{c}\~oes de contorno e a
discretiza\c{c}\~ao do momento.

\paragraph{2. A tend\^encia aponta para o alvo, o que n\~ao \'e o mesmo que
acert\'a-lo.} Em $t=0{,}35$ a velocidade de ascens\~ao est\'a em $0{,}152$ e
subindo; o m\'aximo publicado \'e $0{,}2417$--$0{,}2421$ perto de $t=0{,}92$. A
circularidade come\c{c}a a cair --- $0{,}999974$ para $0{,}999919$ --- como se
espera de uma bolha que come\c{c}a a se deformar. Nenhum desses n\'umeros
\emph{compara} com o publicado, porque nenhum dos tr\^es alvos foi alcan\c{c}ado
no tempo simulado. Registr\'a-los como progresso \'e leg\'itimo; registr\'a-los
como concord\^ancia n\~ao seria.

\subsection{Primeiro alvo fechado: a velocidade m\'axima de ascens\~ao}
\label{sec:vcmax}

\pgfplotstabletypeset[
  string type,
  columns/grandeza/.style={column name={grandeza}},
  columns/FT/.style={column name={front-tracking}},
  columns/VOF/.style={column name={VOF}},
  columns/publicado/.style={column name={publicado}},
  every head row/.style={before row=\toprule, after row=\midrule},
  every last row/.style={after row=\bottomrule},
]{\dat{hysing-vcmax.dat}}

A amostragem \'e a cada $\Delta t = 0{,}05$, e o pico real cai entre amostras;
por isso a linha do refinamento por par\'abola, que corrige \num{4e-5} --- pouco,
o que diz que a amostragem grossa n\~ao \'e o que separa do alvo.

\paragraph{O front-tracking acerta; o VOF n\~ao, por pouco.} O crit\'erio de
compara\c{c}\~ao \'e a dispers\~ao \emph{m\'utua} dos tr\^es c\'odigos publicados:
\SI{0.165}{\percent}. O front-tracking fica a \SI{0.051}{\percent} do limite
inferior da faixa --- \emph{dentro} dela, dado que a faixa \'e a discord\^ancia
entre eles. O VOF fica a \SI{0.378}{\percent}, cerca de $2{,}3$ vezes fora.

Pela leitura declarada na Se\c{c}\~ao~\ref{sec:hysing} \emph{antes} de os
n\'umeros serem conhecidos, este \'e o terceiro desfecho: n\~ao \'e ``os dois
juntos fora da faixa'' --- o que apontaria o que eles compartilham (malha, passo,
contorno, discretiza\c{c}\~ao do momento) --- e sim cada um em sua posi\c{c}\~ao,
com o front-tracking nitidamente mais pr\'oximo. Isso responsabiliza o
\textbf{tratamento de interface}, e \'e coerente com o mecanismo que este
relat\'orio j\'a mediu duas vezes: a curvatura do front-tracking sai da geometria
expl\'icita da frente, a do VOF de derivadas de um campo de fra\c{c}\~ao
suavizado.

Uma ressalva sobre o \emph{instante}: $0{,}9057$ e $0{,}9153$ contra
$0{,}9213$--$0{,}9313$ dos tr\^es c\'odigos. Os dois nossos vem \textbf{cedo}, e o
front-tracking \'e o mais cedo dos cinco. A diferen\c{c}a excede a
dispers\~ao entre eles ($0{,}010$), por\'em pouco.

\subsection{A concord\^ancia m\'utua n\~ao sobreviveu \`a deforma\c{c}\~ao}
\label{sec:divergencia}

\pgfplotstabletypeset[
  string type,
  columns/t/.style={column name={$t$}},
  columns/difyc/.style={column name={dif.\ em $y_c$}},
  columns/difvc/.style={column name={dif.\ em $V_c$}},
  columns/difvcrelativa/.style={column name={dif.\ relativa em $V_c$}},
  every head row/.style={before row=\toprule, after row=\midrule},
  every last row/.style={after row=\bottomrule},
]{\dat{hysing-divergencia.dat}}

Em $t=1{,}6$ a diferen\c{c}a em $V_c$ \'e \SI{3.4}{\percent} do valor. N\~ao \'e
mais ru\'ido num\'erico: \'e discord\^ancia entre os m\'etodos.

\textbf{Isto corrige uma leitura desta pr\'opria se\c{c}\~ao.} A subse\c{c}\~ao do
resultado parcial concluiu, do trecho $t\le0{,}35$, que a concord\^ancia m\'utua
``concentra a suspeita no que \'e comum'' aos dois m\'etodos. A conclus\~ao valia
para o regime \emph{quase circular} em que foi obtida, e n\~ao sobreviveu \`a
deforma\c{c}\~ao --- que \'e exatamente onde este \emph{benchmark} morde. Uma
concord\^ancia medida num regime n\~ao se estende ao seguinte, e aqui n\~ao se
estendeu.

Note a assimetria: $y_c$ estabiliza em $\sim\!\num{4e-5}$ enquanto $V_c$ cresce 68
vezes. \'E o que se espera de uma integral contra a sua derivada --- $y_c$ suaviza
o que $V_c$ amplifica --- e serve de aviso sobre escolher m\'etrica: comparar
apenas centros de massa daria a impress\~ao de concord\^ancia que a velocidade
desmente.

\subsection{Sob malha adaptativa, o VOF perde a conserva\c{c}\~ao exata}
\label{sec:massa-amr}

\pgfplotstabletypeset[
  string type,
  columns/metodo/.style={column name={m\'etodo}},
  columns/malha-fixa/.style={column name={malha fixa}},
  columns/com-AMR-ate-t/.style={column name={com AMR (at\'e $t$)}},
  columns/amostras-monotonas/.style={column name={monotonicidade}},
  every head row/.style={before row=\toprule, after row=\midrule},
  every last row/.style={after row=\bottomrule},
]{\dat{hysing-massa-amr.dat}}

\textbf{Esta \'e a corre\c{c}\~ao mais importante deste relat\'orio.} As
Se\c{c}\~oes~\ref{sec:comparacao}, \ref{sec:balanco} e \ref{sec:elipse} afirmam
que o VOF conserva massa \emph{por constru\c{c}\~ao} --- \num{2.9e-15}, ``zero de
m\'aquina'', ``$A$ em todos os d\'igitos''. Tudo isso foi medido em malha
\textbf{fixa}, e continua verdadeiro ali.

Com remalhamento adaptativo, n\~ao: a fra\c{c}\~ao \'e colhida e replantada a cada
reconstru\c{c}\~ao do dom\'inio, e a deriva chega a \num{1.01e-03} em
$t=2{,}15$ --- doze ordens de grandeza acima do valor em malha fixa, e da
\emph{mesma ordem} da deriva do front-tracking. A monotonicidade descarta ru\'ido:
35 amostras descendo e \textbf{nenhuma} subindo.

A consequ\^encia para a compara\c{c}\~ao \'e direta e desagrad\'avel: a vantagem
que este relat\'orio atribui ao VOF como \emph{estrutural} --- ``conservar \'e
propriedade de constru\c{c}\~ao'' --- \'e estrutural apenas do \emph{esquema de
advec\c{c}\~ao}, e n\~ao do m\'etodo em uso. Basta adaptar a malha para que ela se
dissolva. Num c\'alculo que precise de malha adaptativa, e este precisa, os dois
m\'etodos derivam na mesma ordem.

\subsection{O que decidir\'a}

Quando as corridas fecharem, tr\^es n\'umeros de cada m\'etodo v\~ao contra a
tabela de refer\^encia: $c_{\min}$ e o instante em que ocorre, $V_{c,\max}$ e o
seu instante, e $y_c(t=3)$. A leitura j\'a est\'a decidida de antemão, para que
o resultado n\~ao seja interpretado depois de conhecido:

\begin{itemize}
  \item \textbf{Dentro da faixa de consenso} nos tr\^es: o par de
    implementa\c{c}\~oes est\'a verificado contra a literatura, e a compara\c{c}\~ao
    FT $\times$ VOF deste relat\'orio passa a valer como compara\c{c}\~ao entre
    m\'etodos, e n\~ao apenas entre duas implementa\c{c}\~oes nossas.
  \item \textbf{Fora da faixa, e os dois m\'etodos juntos}: o defeito est\'a no
    que eles compartilham --- malha, passo, contorno, discretiza\c{c}\~ao do
    momento --- e n\~ao no tratamento de interface. \'E o desfecho mais
    prov\'avel dado o padr\~ao da tabela parcial, e o mais informativo, porque
    aponta um lugar estreito.
  \item \textbf{Fora da faixa, e cada um para um lado}: o tratamento de interface
    responde, e a diferen\c{c}a entre as duas colunas diz qual.
\end{itemize}

As tr\^es leituras s\~ao distintas e caem em lugares diferentes do c\'odigo;
declar\'a-las antes \'e o que impede a quarta, que \'e explicar qualquer
resultado como esperado.
